\begin{ouverture}[statistiques]{18}{Organiser et interpréter des données}{Partie 4 · Données, hasard et programmation}
{Effectifs · Fréquences · Diagrammes · Moyenne simple}
\objectif{Recueillir et organiser des données.}
\objectif{Calculer des fréquences et choisir une représentation.}
\objectif{Calculer et interpréter une moyenne simple.}
\prerequis{Fractions, pourcentages et graphiques.}
\end{ouverture}
\pagestyle{otMain}\essentiel\fontsize{11}{15}\selectfont
\phantomsection\label{cours-18-1}
\essSection{1}{Du relevé à la représentation}
\begin{definition}Une série statistique rassemble des données. L’effectif d’une valeur ou d’une catégorie est son nombre d’apparitions. L’effectif total est la somme des effectifs. La fréquence d’une catégorie est son effectif divisé par le total non nul.\end{definition}
\begin{center}\renewcommand{\arraystretch}{1.35}\begin{tabular}{c|ccc}\toprule
Transport & À pied & Vélo & Bus\\\midrule
Effectif & 10 & 6 & 4\\
Fréquence & 50 \% & 30 \% & 20 \%\\
\bottomrule\end{tabular}\end{center}
\label{fig-cours-18-1-cours-statistiques}\figureGeo{cours-statistiques}

\begin{savoirfaire}[Choisir pour montrer]Un diagramme en barres compare les catégories ; un diagramme circulaire montre le partage du total ; un graphique cartésien montre une évolution. Choisir la représentation selon la question, avec titres, unités et graduations.\end{savoirfaire}
\clearpage
\phantomsection\label{cours-18-2}
\essSection{2}{Moyenne et lecture critique}
\begin{definition}La moyenne simple d’une série numérique est la somme des valeurs divisée par leur nombre. Elle décrit une valeur de partage égal ; elle n’est pas nécessairement une valeur observée.\end{definition}
\begin{exemple}Pour $4,6,6,8$, la somme est 24 et il y a quatre données : la moyenne vaut $24:4=6$. Pour $2,3,10$, elle vaut 5, même si aucun résultat ne vaut 5. Si l’on ajoute 1 à chacune des données $4,6,6,8$, la somme augmente de 4 et la moyenne devient $28:4=7$ : elle augmente aussi de 1.\end{exemple}
\begin{savoirfaire}[Exploiter un tableur]Dans un tableur, placer une donnée par cellule ; \texttt{=MOYENNE(A2:A6)} calcule la moyenne de cinq données. Pour une fréquence, diviser l’effectif par le total puis choisir l’affichage en pourcentage.\end{savoirfaire}
\begin{attention}Une moyenne seule ne montre pas la dispersion des valeurs. Deux séries peuvent avoir la même moyenne et des profils très différents. La hauteur d’une barre ne se lit qu’après vérification de l’échelle.\end{attention}

\clearpage\rubrique{Découverte}
\begin{blocLivre}{ex-18-1}\begin{exercice}[1]{Un effectif total}
Un relevé comporte 7 réponses A, 9 réponses B et 4 réponses C. Donner l’effectif total.
\end{exercice}\end{blocLivre}
\begin{blocLivre}{ex-18-2}\begin{exercice}[1]{Une fréquence}
Sur 25 élèves, 10 choisissent le vélo. Donner la fréquence sous forme de fraction, de décimal et de pourcentage.
\end{exercice}\end{blocLivre}
\begin{blocLivre}{ex-18-3}\begin{exercice}[1]{Une moyenne directe}
Calculer la moyenne de 8, 10, 12 et 14.
\end{exercice}\end{blocLivre}
\begin{blocLivre}{ex-18-4}\begin{exercice}[1]{Lire des barres}
Le diagramme donne le nombre de livres lus. Combien d’élèves ont lu deux livres ?
\label{fig-ex-18-4}\figureGeo{stat-4}
\end{exercice}\end{blocLivre}
\begin{blocLivre}{ex-18-5}\begin{exercice}[1]{Choisir un graphique}
Pour comparer les effectifs de trois sports et pour montrer une température au fil des heures, choisir deux représentations adaptées.
\end{exercice}\end{blocLivre}

\rubrique{Entraînement}
\begin{blocLivre}{ex-18-6}\begin{exercice}[2]{Organiser une liste}
On observe les valeurs $1,2,1,3,2,1,2,2$. Donner les effectifs et fréquences de 1, 2 et 3.
\end{exercice}\end{blocLivre}
\begin{blocLivre}{ex-18-7}\begin{exercice}[2]{Construire des secteurs}
Sur 20 votes, A en obtient 10, B en obtient 6 et C en obtient 4. Donner les angles du diagramme circulaire puis le construire.
\label{fig-ex-18-7}\figureGeo{stat-7}
\end{exercice}\end{blocLivre}
\begin{blocLivre}{ex-18-8}\begin{exercice}[2]{Lire un secteur}
Un secteur de diagramme circulaire mesure $90^\circ$ sur un total de 48 personnes. Combien représente-t-il de personnes ?
\label{fig-ex-18-8}\figureGeo{stat-8}
\end{exercice}\end{blocLivre}
\begin{blocLivre}{ex-18-9}\begin{exercice}[2]{Une moyenne décimale}
Quatre mesures valent 2,4 ; 3,1 ; 2,8 et 3,7 cm. Calculer la moyenne. On ajoute ensuite 0,5 cm à chaque mesure : quelle est la nouvelle moyenne ?
\end{exercice}\end{blocLivre}
\begin{blocLivre}{ex-18-10}\begin{exercice}[2]{Une donnée extrême}
Calculer la moyenne de 4, 5, 6 ; puis celle de 4, 5, 60. Quel effet produit la dernière valeur ?
\end{exercice}\end{blocLivre}

\rubrique{Problèmes}
\begin{blocLivre}{ex-18-11}\begin{exercice}[3]{Retrouver une donnée}
La moyenne de cinq nombres est 12. Quatre de ces nombres sont 8, 11, 13 et 15. Trouver le cinquième.
\end{exercice}\end{blocLivre}
\begin{blocLivre}{ex-18-12}\begin{exercice}[3]{Comparer des groupes}
Dans A, 12 personnes sur 20 ont réussi ; dans B, 18 sur 40. Quel groupe a le plus grand nombre de réussites ? Le meilleur taux ?
\end{exercice}\end{blocLivre}
\begin{blocLivre}{ex-18-13}\begin{exercice}[3]{Même moyenne, profils différents}
Comparer les séries $5,5,5,5$ et $0,0,10,10$. Calculer leurs moyennes et expliquer ce que la moyenne ne dit pas.
\end{exercice}\end{blocLivre}
\begin{blocLivre}{ex-18-14}\begin{exercice}[3]{Une graduation trompeuse}
Le diagramme compare 48 et 52 personnes avec un axe vertical commençant à 40. Les barres semblent avoir des hauteurs 8 et 12. Le second effectif est-il 1,5 fois le premier ?
\label{fig-ex-18-14}\figureGeo{stat-14}
\end{exercice}\end{blocLivre}
\begin{blocLivre}{ex-18-15}\begin{exercice}[3]{Préparer une enquête}
On veut connaître le moyen de transport principal de 30 élèves, une réponse par élève. Les réponses donnent 12 à pied, 9 en bus, 6 à vélo et 3 autrement. Donner les fréquences et proposer deux graphiques possibles.
\label{fig-ex-18-15}\figureGeo{stat-15}
\end{exercice}\end{blocLivre}

\clearpage
